% ============================================================
\chapter{Stabilit\'e de Lyapunov}
\label{ch:lyapunov}
% ============================================================

En 1892, le math\'ematicien russe Aleksandr Lyapunov soutient sa th\`ese
\emph{Le probl\`eme g\'en\'eral de la stabilit\'e du mouvement}. Son id\'ee
g\'eniale~: pour d\'ecider si un point d'\'equilibre est stable, il n'est pas
n\'ecessaire de r\'esoudre les \'equations du mouvement. Il suffit de trouver
une \og fonction d'\'energie\fg{} --- aujourd'hui appel\'ee \emph{fonction de
Lyapunov} --- qui d\'ecro\^it le long des trajectoires. Si une telle
fonction existe, le syst\`eme est stable, \emph{sans qu'on ait besoin de
conna\^itre les solutions explicitement}. C'est l'une des m\'ethodes les plus
puissantes de toute la th\'eorie des syst\`emes dynamiques.

% ------------------------------------------------------------
\section{Introduction}
% ------------------------------------------------------------

La question de la stabilit\'e est fondamentale dans l'\'etude des syst\`emes
dynamiques : un point d'\'equilibre est-il robuste face aux perturbations ?
Alexandre Lyapounov a d\'evelopp\'e dans sa th\`ese de 1892 deux approches
compl\'ementaires : la m\'ethode indirecte (lin\'earisation) et la m\'ethode
directe (fonctions de Lyapounov). Ces m\'ethodes restent les outils centraux
de l'analyse de stabilit\'e.

\section{D\'efinitions de la stabilit\'e}

Consid\'erons le syst\`eme autonome $\dot{x} = f(x)$ o\`u $f : U \to \R^n$
est de classe $C^1$ et $x^*$ est un point d'\'equilibre : $f(x^*) = 0$.
Sans perte de g\'en\'eralit\'e, nous supposons $x^* = 0$ (par translation).

\begin{definition}[Stabilit\'e au sens de Lyapunov]
Le point d'\'equilibre $x^* = 0$ est \textbf{stable au sens de Lyapunov} si
pour tout $\epsilon > 0$, il existe $\delta > 0$ tel que
\[
    \norm{x_0} < \delta \implies \norm{\varphi(t, x_0)} < \epsilon
    \quad \forall t \geq 0.
\]
\end{definition}

\begin{definition}[Stabilit\'e asymptotique]
Le point d'\'equilibre $0$ est \textbf{asymptotiquement stable} s'il est
stable au sens de Lyapunov et s'il existe $\delta_0 > 0$ tel que
\[
    \norm{x_0} < \delta_0 \implies \lim_{t \to +\infty} \varphi(t, x_0) = 0.
\]
\end{definition}

\begin{definition}[Stabilit\'e asymptotique globale]
Le point $0$ est \textbf{globalement asymptotiquement stable} (GAS) s'il est
stable et si $\lim_{t \to +\infty} \varphi(t, x_0) = 0$ pour tout $x_0 \in U$.
\end{definition}

\begin{definition}[Instabilit\'e]
Le point $0$ est \textbf{instable} s'il n'est pas stable au sens de Lyapunov.
\end{definition}

\begin{definition}[Stabilit\'e exponentielle]
Le point $0$ est \textbf{exponentiellement stable} s'il existe des constantes
$M > 0$, $\alpha > 0$ et $\delta > 0$ tels que
\[
    \norm{x_0} < \delta \implies \norm{\varphi(t, x_0)} \leq M\norm{x_0}
    e^{-\alpha t} \quad \forall t \geq 0.
\]
\end{definition}

\begin{remarque}
La stabilit\'e exponentielle implique la stabilit\'e asymptotique, qui implique
la stabilit\'e au sens de Lyapunov. Les r\'eciproques sont fausses en g\'en\'eral.
\end{remarque}

\begin{intuition}
\begin{itemize}
    \item \textbf{Stabilit\'e de Lyapunov} : une bille au fond d'un bol reste
    pr\`es du fond si on la perturbe l\'eg\`erement.
    \item \textbf{Stabilit\'e asymptotique} : la bille revient au fond gr\^ace
    au frottement.
    \item \textbf{Instabilit\'e} : une bille au sommet d'une colline ---
    toute perturbation l'\'eloigne.
\end{itemize}
\end{intuition}

\section{Bassin d'attraction}

\begin{definition}[Bassin d'attraction]
Le \textbf{bassin d'attraction} d'un point d'\'equilibre asymptotiquement
stable $x^*$ est l'ensemble
\[
    \mathcal{B}(x^*) = \{x_0 \in U : \lim_{t \to +\infty} \varphi(t, x_0) = x^*\}.
\]
\end{definition}

\begin{proposition}
Le bassin d'attraction $\mathcal{B}(x^*)$ est un ouvert invariant par le flot,
et sa fronti\`ere $\partial\mathcal{B}(x^*)$ est \'egalement invariante.
\end{proposition}

\section{M\'ethode indirecte : lin\'earisation}

\begin{theoreme}[Stabilit\'e par lin\'earisation --- Lyapunov indirect]
\label{thm:lyap-indirect}
Soit $A = Df(0)$ la matrice jacobienne de $f$ au point d'\'equilibre $x^* = 0$.
\begin{enumerate}
    \item Si toutes les valeurs propres de $A$ ont une partie r\'eelle
    strictement n\'egative ($\text{Re}(\lambda_i) < 0$ pour tout $i$),
    alors $0$ est \textbf{asymptotiquement stable}.
    \item Si au moins une valeur propre de $A$ a une partie r\'eelle
    strictement positive ($\text{Re}(\lambda_j) > 0$ pour un $j$),
    alors $0$ est \textbf{instable}.
    \item Si toutes les valeurs propres satisfont $\text{Re}(\lambda_i) \leq 0$
    avec au moins une valeur propre sur l'axe imaginaire, la lin\'earisation
    \textbf{ne conclut pas} : il faut utiliser la m\'ethode directe.
\end{enumerate}
\end{theoreme}

\begin{proof}[Id\'ee de la d\'emonstration]
On \'ecrit $f(x) = Ax + g(x)$ avec $g(x) = o(\norm{x})$ au voisinage de $0$.
Si toutes les valeurs propres de $A$ ont des parties r\'eelles n\'egatives,
il existe $P$ d\'efinie positive telle que $A^\top P + PA = -Q$ avec $Q$
d\'efinie positive (\'equation de Lyapounov matricielle). La fonction
$V(x) = x^\top P x$ est alors une fonction de Lyapounov pour le syst\`eme
lin\'earis\'e, et le terme $g(x) = o(\norm{x})$ ne perturbe pas la conclusion
dans un voisinage suffisamment petit de l'origine.
\end{proof}

\begin{exemple}[Pendule amorti]
Le pendule amorti lin\'earis\'e au voisinage de $\theta = 0$ donne
\[
    A = \begin{pmatrix} 0 & 1 \\ -g/l & -c \end{pmatrix},
\]
avec $c > 0$ le coefficient d'amortissement. Le polyn\^ome caract\'eristique est
$\lambda^2 + c\lambda + g/l = 0$, de discriminant $\Delta = c^2 - 4g/l$.
Les deux valeurs propres ont des parties r\'eelles n\'egatives
(car $c > 0$ et $g/l > 0$), donc l'\'equilibre est asymptotiquement stable.
\end{exemple}

\section{M\'ethode directe de Lyapunov}

\begin{definition}[Fonction de Lyapunov]
Soit $V : U \to \R$ une fonction de classe $C^1$ d\'efinie sur un voisinage
de $x^* = 0$ avec $V(0) = 0$. La \textbf{d\'eriv\'ee orbitale} (ou d\'eriv\'ee
de Lie) de $V$ le long du champ $f$ est
\[
    \dot{V}(x) = \nabla V(x) \cdot f(x) = \sum_{i=1}^n \frac{\partial V}{\partial x_i}(x)\, f_i(x).
\]
On dit que $V$ est :
\begin{itemize}
    \item \textbf{d\'efinie positive} si $V(x) > 0$ pour tout $x \neq 0$
    dans un voisinage de $0$ ;
    \item \textbf{semi-d\'efinie positive} si $V(x) \geq 0$ ;
    \item \textbf{d\'efinie n\'egative} si $V(x) < 0$ pour tout $x \neq 0$.
\end{itemize}
\end{definition}

\begin{theoreme}[Th\'eor\`eme de stabilit\'e de Lyapunov]
\label{thm:lyap-stable}
S'il existe une fonction $V$ d\'efinie positive telle que $\dot{V}(x) \leq 0$
dans un voisinage de $0$, alors $0$ est \textbf{stable au sens de Lyapunov}.
\end{theoreme}

\begin{proof}
Soit $\epsilon > 0$ suffisamment petit pour que $\overline{B}(0, \epsilon) \subseteq U$.
Posons $c = \min_{\norm{x} = \epsilon} V(x) > 0$ (car $V$ est d\'efinie positive
et continue sur le compact $\{\norm{x} = \epsilon\}$). Par continuit\'e de $V$
et $V(0) = 0$, il existe $\delta > 0$ tel que $V(x) < c$ pour $\norm{x} < \delta$.
Si $\norm{x_0} < \delta$, alors $V(\varphi(t, x_0)) \leq V(x_0) < c$ pour tout
$t \geq 0$ (car $\dot{V} \leq 0$). Donc $\varphi(t, x_0)$ ne peut pas atteindre
la sph\`ere $\norm{x} = \epsilon$, ce qui prouve $\norm{\varphi(t, x_0)} < \epsilon$.
\end{proof}

\begin{theoreme}[Th\'eor\`eme de stabilit\'e asymptotique de Lyapunov]
\label{thm:lyap-asymptotic}
S'il existe une fonction $V$ d\'efinie positive telle que $\dot{V}(x) < 0$
pour tout $x \neq 0$ dans un voisinage de $0$ (i.e., $\dot{V}$ est
d\'efinie n\'egative), alors $0$ est \textbf{asymptotiquement stable}.
\end{theoreme}

\begin{theoreme}[Th\'eor\`eme de Barbashin-Krasovskii]
\label{thm:barbashin}
Si $V : \R^n \to \R$ est d\'efinie positive, radialement non born\'ee
($V(x) \to +\infty$ quand $\norm{x} \to +\infty$), et si $\dot{V}(x) < 0$
pour tout $x \neq 0$, alors $0$ est \textbf{globalement asymptotiquement stable}.
\end{theoreme}

\begin{attention}[title=\textbf{Fonctions de Lyapunov non born\'ees}]
Pour la stabilit\'e asymptotique \emph{globale}, il est essentiel que $V$ soit
radialement non born\'ee. Sans cette condition, la trajectoire pourrait
s'\'echapper vers l'infini tout en gardant $V$ d\'ecroissante.
\end{attention}

\section{Principe d'invariance de LaSalle}

Lorsque $\dot{V}$ est seulement semi-d\'efinie n\'egative, le th\'eor\`eme de
stabilit\'e asymptotique ne s'applique pas directement. Le principe de LaSalle
permet souvent de conclure malgr\'e tout.

\begin{theoreme}[Principe d'invariance de LaSalle]
\label{thm:lasalle}
Soit $\Omega \subseteq U$ un compact positivement invariant. Soit
$V : \Omega \to \R$ de classe $C^1$ telle que $\dot{V}(x) \leq 0$ sur $\Omega$.
Posons $S = \{x \in \Omega : \dot{V}(x) = 0\}$ et soit $M$ le plus grand
sous-ensemble invariant contenu dans $S$. Alors
\[
    \omega(x_0) \subseteq M \qquad \forall x_0 \in \Omega.
\]
En particulier, toute trajectoire issue de $\Omega$ converge vers $M$.
\end{theoreme}

\begin{proof}[Id\'ee de la d\'emonstration]
Puisque $V$ est d\'ecroissante et born\'ee inf\'erieurement sur le compact $\Omega$,
$V(\varphi(t, x_0))$ converge vers une limite $c$. L'ensemble $\omega$-limite
$\omega(x_0)$ est non vide (car $\Omega$ est compact), invariant, et
$V \equiv c$ sur $\omega(x_0)$. Donc $\dot{V} = 0$ sur $\omega(x_0)$,
ce qui donne $\omega(x_0) \subseteq S$. Par invariance de $\omega(x_0)$,
on obtient $\omega(x_0) \subseteq M$.
\end{proof}

\begin{exemple}[Pendule amorti --- LaSalle]
Consid\'erons le pendule amorti :
\[
    \dot{\theta} = \omega, \qquad \dot{\omega} = -\frac{g}{l}\sin\theta - c\omega,
    \quad c > 0.
\]
La fonction \'energie $V(\theta, \omega) = \frac{1}{2}\omega^2 + \frac{g}{l}(1 - \cos\theta)$
satisfait
\[
    \dot{V} = -c\omega^2 \leq 0.
\]
L'ensemble $S = \{\dot{V} = 0\} = \{\omega = 0\}$. Sur $S$, on a
$\dot{\omega} = -\frac{g}{l}\sin\theta$, qui est nul uniquement si
$\theta = k\pi$. Le plus grand ensemble invariant dans $S$ est
$M = \{(0, 0), (\pi, 0)\}$. Par LaSalle, toute trajectoire born\'ee converge
vers $M$. Pr\`es de l'\'equilibre $(0,0)$, on obtient la stabilit\'e asymptotique.
\end{exemple}

\section{Th\'eor\`emes d'instabilit\'e}

\begin{theoreme}[Th\'eor\`eme d'instabilit\'e de Lyapunov]
\label{thm:instabilite-lyap}
S'il existe une fonction $V$ de classe $C^1$ avec $V(0) = 0$ telle que
$\dot{V}(x) > 0$ dans un voisinage de $0$ (sauf en $0$), et si dans tout
voisinage de $0$ il existe un point $x$ avec $V(x) > 0$, alors $0$ est
\textbf{instable}.
\end{theoreme}

\begin{theoreme}[Th\'eor\`eme de Chetaev]
Soit $V : U \to \R$ de classe $C^1$ avec $V(0) = 0$. Supposons qu'il existe
un ouvert $G$ tel que $0 \in \partial G$ et :
\begin{enumerate}
    \item $V(x) > 0$ pour tout $x \in G$ ;
    \item $V(x) = 0$ sur $\partial G \cap U$ ;
    \item $\dot{V}(x) > 0$ pour tout $x \in G$.
\end{enumerate}
Alors $0$ est \textbf{instable}.
\end{theoreme}

\section{Construction de fonctions de Lyapunov}

La difficult\'e principale de la m\'ethode directe est de \emph{trouver} une
fonction de Lyapunov ad\'equate. Voici quelques strat\'egies.

\subsection{Fonctions quadratiques}

Pour le syst\`eme $\dot{x} = Ax + g(x)$ avec $g(x) = o(\norm{x})$, on
cherche $V(x) = x^\top P x$ avec $P$ sym\'etrique d\'efinie positive solution
de l'\'equation de Lyapounov matricielle :
\[
    A^\top P + PA = -Q,
\]
o\`u $Q$ est sym\'etrique d\'efinie positive. Cette \'equation admet une unique
solution $P$ d\'efinie positive si et seulement si $A$ est une matrice de
Hurwitz (toutes les valeurs propres \`a partie r\'eelle n\'egative).

\begin{formules}[title=\textbf{\'Equation de Lyapounov matricielle}]
\[
    A^\top P + PA = -Q \qquad \Longrightarrow \qquad
    P = \int_0^{+\infty} e^{A^\top t}\, Q\, e^{At}\, dt.
\]
Cette int\'egrale converge si et seulement si $A$ est une matrice de Hurwitz.
\end{formules}

\subsection{Fonctions \'energie}

Pour les syst\`emes m\'ecaniques, l'\'energie totale (cin\'etique + potentielle)
est souvent un bon candidat. Si le syst\`eme est dissipatif, $\dot{V}$ sera
n\'egative ou semi-n\'egative.

\subsection{Combinaisons et estimations}

Pour des syst\`emes plus complexes, on peut construire $V$ par combinaison
lin\'eaire de termes simples, ou par la m\'ethode des \emph{backstepping} en
th\'eorie du contr\^ole.

\section{Stabilit\'e des orbites p\'eriodiques}

\begin{definition}[Stabilit\'e orbitale]
Une orbite p\'eriodique $\Gamma$ est \textbf{orbitalement stable} si pour tout
$\epsilon > 0$, il existe $\delta > 0$ tel que
\[
    \mathrm{dist}(x_0, \Gamma) < \delta \implies
    \mathrm{dist}(\varphi(t, x_0), \Gamma) < \epsilon \quad \forall t \geq 0.
\]
\end{definition}

\begin{definition}[Stabilit\'e asymptotique orbitale]
L'orbite $\Gamma$ est \textbf{asymptotiquement orbitalement stable} si elle
est orbitalement stable et si
$\mathrm{dist}(\varphi(t, x_0), \Gamma) \to 0$ quand $t \to +\infty$
pour $x_0$ suffisamment proche de $\Gamma$.
\end{definition}

\begin{remarque}
La stabilit\'e orbitale est plus naturelle que la stabilit\'e au sens de Lyapunov
pour les orbites p\'eriodiques, car m\^eme dans le cas stable, deux points
voisins sur l'orbite se s\'eparent lin\'eairement en temps le long de l'orbite.
\end{remarque}

\section{Fonctions de Lyapunov et contr\^ole}

\begin{remarque}
En th\'eorie du contr\^ole, les fonctions de Lyapunov jouent un r\^ole central.
Pour un syst\`eme contr\^ol\'e $\dot{x} = f(x, u)$, on cherche un retour
d'\'etat $u = k(x)$ tel qu'une fonction de Lyapunov $V$ satisfasse
$\dot{V} < 0$ en boucle ferm\'ee. C'est le principe du
\textbf{contr\^ole par fonctions de Lyapunov} (CLF --- Control Lyapunov Function).
\end{remarque}

\section{Exercices}

\begin{exercice}[Stabilit\'e par lin\'earisation]
\'Etudier la stabilit\'e des points d'\'equilibre du syst\`eme
\[
    \dot{x} = y + x(1 - x^2 - y^2), \qquad \dot{y} = -x + y(1 - x^2 - y^2).
\]
\end{exercice}

\begin{exercice}[Construction de Lyapunov]
Pour le syst\`eme $\dot{x} = -x + 2x^2y$, $\dot{y} = -y$, montrer que
$V(x,y) = x^2 + y^2$ est une fonction de Lyapunov pour un voisinage
convenable de l'origine et conclure sur la stabilit\'e.
\end{exercice}

\begin{exercice}[LaSalle]
Appliquer le principe de LaSalle au syst\`eme
\[
    \dot{x} = y, \qquad \dot{y} = -x^3 - y^3,
\]
avec la fonction $V(x,y) = \frac{x^4}{4} + \frac{y^2}{2}$.
\end{exercice}

\begin{exercice}[Instabilit\'e par Chetaev]
Montrer que l'origine est instable pour le syst\`eme
\[
    \dot{x} = y + x^3, \qquad \dot{y} = x + y^3,
\]
en utilisant la fonction $V(x,y) = xy$.
\end{exercice}

\begin{exercice}[\'Equation de Lyapounov matricielle]
R\'esoudre l'\'equation $A^\top P + PA = -I$ pour
$A = \begin{pmatrix} -1 & 0 \\ 1 & -2 \end{pmatrix}$.
V\'erifier que $P$ est d\'efinie positive.
\end{exercice}

\begin{exercice}[Stabilit\'e globale]
Montrer que l'origine est globalement asymptotiquement stable pour
$\dot{x} = -x^3$, $\dot{y} = -y$ en utilisant $V(x,y) = \frac{x^4}{4} + \frac{y^2}{2}$.
\end{exercice}
